% Options for packages loaded elsewhere
\PassOptionsToPackage{unicode}{hyperref}
\PassOptionsToPackage{hyphens}{url}
\documentclass[
  ignorenonframetext,
]{beamer}
\newif\ifbibliography
\usepackage{pgfpages}
\setbeamertemplate{caption}[numbered]
\setbeamertemplate{caption label separator}{: }
\setbeamercolor{caption name}{fg=normal text.fg}
\beamertemplatenavigationsymbolsempty
% remove section numbering
\setbeamertemplate{part page}{
  \centering
  \begin{beamercolorbox}[sep=16pt,center]{part title}
    \usebeamerfont{part title}\insertpart\par
  \end{beamercolorbox}
}
\setbeamertemplate{section page}{
  \centering
  \begin{beamercolorbox}[sep=12pt,center]{section title}
    \usebeamerfont{section title}\insertsection\par
  \end{beamercolorbox}
}
\setbeamertemplate{subsection page}{
  \centering
  \begin{beamercolorbox}[sep=8pt,center]{subsection title}
    \usebeamerfont{subsection title}\insertsubsection\par
  \end{beamercolorbox}
}
% Prevent slide breaks in the middle of a paragraph
\widowpenalties 1 10000
\raggedbottom
\AtBeginPart{
  \frame{\partpage}
}
\AtBeginSection{
  \ifbibliography
  \else
    \frame{\sectionpage}
  \fi
}
\AtBeginSubsection{
  \frame{\subsectionpage}
}
\usepackage{iftex}
\ifPDFTeX
  \usepackage[T1]{fontenc}
  \usepackage[utf8]{inputenc}
  \usepackage{textcomp} % provide euro and other symbols
\else % if luatex or xetex
  \usepackage{unicode-math} % this also loads fontspec
  \defaultfontfeatures{Scale=MatchLowercase}
  \defaultfontfeatures[\rmfamily]{Ligatures=TeX,Scale=1}
\fi
\usepackage{lmodern}
\ifPDFTeX\else
  % xetex/luatex font selection
\fi
% Use upquote if available, for straight quotes in verbatim environments
\IfFileExists{upquote.sty}{\usepackage{upquote}}{}
\IfFileExists{microtype.sty}{% use microtype if available
  \usepackage[]{microtype}
  \UseMicrotypeSet[protrusion]{basicmath} % disable protrusion for tt fonts
}{}
\makeatletter
\@ifundefined{KOMAClassName}{% if non-KOMA class
  \IfFileExists{parskip.sty}{%
    \usepackage{parskip}
  }{% else
    \setlength{\parindent}{0pt}
    \setlength{\parskip}{6pt plus 2pt minus 1pt}}
}{% if KOMA class
  \KOMAoptions{parskip=half}}
\makeatother
\usepackage{color}
\usepackage{fancyvrb}
\newcommand{\VerbBar}{|}
\newcommand{\VERB}{\Verb[commandchars=\\\{\}]}
\DefineVerbatimEnvironment{Highlighting}{Verbatim}{commandchars=\\\{\}}
% Add ',fontsize=\small' for more characters per line
\newenvironment{Shaded}{}{}
\newcommand{\AlertTok}[1]{\textcolor[rgb]{1.00,0.00,0.00}{\textbf{#1}}}
\newcommand{\AnnotationTok}[1]{\textcolor[rgb]{0.38,0.63,0.69}{\textbf{\textit{#1}}}}
\newcommand{\AttributeTok}[1]{\textcolor[rgb]{0.49,0.56,0.16}{#1}}
\newcommand{\BaseNTok}[1]{\textcolor[rgb]{0.25,0.63,0.44}{#1}}
\newcommand{\BuiltInTok}[1]{\textcolor[rgb]{0.00,0.50,0.00}{#1}}
\newcommand{\CharTok}[1]{\textcolor[rgb]{0.25,0.44,0.63}{#1}}
\newcommand{\CommentTok}[1]{\textcolor[rgb]{0.38,0.63,0.69}{\textit{#1}}}
\newcommand{\CommentVarTok}[1]{\textcolor[rgb]{0.38,0.63,0.69}{\textbf{\textit{#1}}}}
\newcommand{\ConstantTok}[1]{\textcolor[rgb]{0.53,0.00,0.00}{#1}}
\newcommand{\ControlFlowTok}[1]{\textcolor[rgb]{0.00,0.44,0.13}{\textbf{#1}}}
\newcommand{\DataTypeTok}[1]{\textcolor[rgb]{0.56,0.13,0.00}{#1}}
\newcommand{\DecValTok}[1]{\textcolor[rgb]{0.25,0.63,0.44}{#1}}
\newcommand{\DocumentationTok}[1]{\textcolor[rgb]{0.73,0.13,0.13}{\textit{#1}}}
\newcommand{\ErrorTok}[1]{\textcolor[rgb]{1.00,0.00,0.00}{\textbf{#1}}}
\newcommand{\ExtensionTok}[1]{#1}
\newcommand{\FloatTok}[1]{\textcolor[rgb]{0.25,0.63,0.44}{#1}}
\newcommand{\FunctionTok}[1]{\textcolor[rgb]{0.02,0.16,0.49}{#1}}
\newcommand{\ImportTok}[1]{\textcolor[rgb]{0.00,0.50,0.00}{\textbf{#1}}}
\newcommand{\InformationTok}[1]{\textcolor[rgb]{0.38,0.63,0.69}{\textbf{\textit{#1}}}}
\newcommand{\KeywordTok}[1]{\textcolor[rgb]{0.00,0.44,0.13}{\textbf{#1}}}
\newcommand{\NormalTok}[1]{#1}
\newcommand{\OperatorTok}[1]{\textcolor[rgb]{0.40,0.40,0.40}{#1}}
\newcommand{\OtherTok}[1]{\textcolor[rgb]{0.00,0.44,0.13}{#1}}
\newcommand{\PreprocessorTok}[1]{\textcolor[rgb]{0.74,0.48,0.00}{#1}}
\newcommand{\RegionMarkerTok}[1]{#1}
\newcommand{\SpecialCharTok}[1]{\textcolor[rgb]{0.25,0.44,0.63}{#1}}
\newcommand{\SpecialStringTok}[1]{\textcolor[rgb]{0.73,0.40,0.53}{#1}}
\newcommand{\StringTok}[1]{\textcolor[rgb]{0.25,0.44,0.63}{#1}}
\newcommand{\VariableTok}[1]{\textcolor[rgb]{0.10,0.09,0.49}{#1}}
\newcommand{\VerbatimStringTok}[1]{\textcolor[rgb]{0.25,0.44,0.63}{#1}}
\newcommand{\WarningTok}[1]{\textcolor[rgb]{0.38,0.63,0.69}{\textbf{\textit{#1}}}}
\usepackage{longtable,booktabs,array}
\newcounter{none} % for unnumbered tables
\usepackage{calc} % for calculating minipage widths
\usepackage{caption}
% Make caption package work with longtable
\makeatletter
\def\fnum@table{\tablename~\thetable}
\makeatother
\setlength{\emergencystretch}{3em} % prevent overfull lines
\providecommand{\tightlist}{%
  \setlength{\itemsep}{0pt}\setlength{\parskip}{0pt}}
% Slide layout defaults for warning-clean scientific decks.
\usepackage{etoolbox}
\IfFileExists{xurl.sty}{\usepackage{xurl}}{}
\IfFileExists{seqsplit.sty}{\usepackage{seqsplit}}{\newcommand{\seqsplit}[1]{#1}}
\protected\def\breakseq#1{\seqsplit{#1}}
\protected\def\breaktt#1{\begingroup\ttfamily\seqsplit{#1}\endgroup}
\setlength{\emergencystretch}{6em}
\tolerance=5000
\hbadness=10000
\hfuzz=1pt
\setlength{\tabcolsep}{2pt}
\AtBeginEnvironment{longtable}{\tiny\renewcommand{\arraystretch}{0.86}\setlength{\tabcolsep}{1pt}}
\AtBeginEnvironment{tabular}{\tiny\renewcommand{\arraystretch}{0.86}\setlength{\tabcolsep}{1pt}}

% Natbib and cross-reference fallbacks — slides don't load natbib
% or cleveref, but combined-PDF manuscript prose may emit these
% commands. The fallback renders citations as a bracketed key list
% and cross-references as detokenized labels so slides don't fail on
% undefined control sequences or raw underscores. \providecommand is
% a no-op if packages load later.
\providecommand{\citep}[1]{[#1]}
\providecommand{\citet}[1]{#1}
\providecommand{\citealp}[1]{#1}
\providecommand{\citeauthor}[1]{#1}
\providecommand{\citeyear}[1]{#1}
\providecommand{\cref}[1]{\texttt{\detokenize{#1}}}
\providecommand{\Cref}[1]{\texttt{\detokenize{#1}}}
\usepackage{bookmark}
\IfFileExists{xurl.sty}{\usepackage{xurl}}{} % add URL line breaks if available
\urlstyle{same}
% fallback for those not using the hyperref driver hyperxmp:
\makeatletter
\@ifundefined{xmpquote}{\newcommand{\xmpquote}[1]{#1}}{}
\makeatother
\hypersetup{
  hidelinks,
  pdfcreator={LaTeX via pandoc}}

\author{\texorpdfstring{}{}}
\date{}

\begin{document}

\begin{frame}[fragile,allowframebreaks]{Documentation Duality and AI
Collaboration}
\protect\phantomsection\label{documentation-duality-and-ai-collaboration}
Every directory at every level of the repository hierarchy contains two
documentation files:

\begin{itemize}
\tightlist
\item
  \textbf{\texttt{README.md}}: Human-readable overview, quick-start
  instructions, and directory structure.
\item
  \textbf{\texttt{AGENTS.md}}: Machine-readable technical specification
  optimized for AI coding assistants. Contains API tables, dependency
  graphs, implementation patterns, and architectural constraints.
\end{itemize}

This Documentation Duality standard serves two purposes. First, it
ensures that both human researchers and AI agents can navigate the
codebase efficiently---\texttt{AGENTS.md} files provide the structured
context that language models need to make informed code modifications
without hallucinating API signatures or violating architectural
invariants. Second, it creates a self-documenting feedback loop: as AI
agents modify the codebase, they update the corresponding
\texttt{AGENTS.md} files, keeping documentation synchronized with
implementation. Lau and Guo's survey of 90 AI coding assistant systems
\breakseq{[@lau2025aicoding]} identifies contextual code understanding as a
primary bottleneck; the Documentation Duality standard addresses this by
providing pre-structured context at every directory level.

The template additionally includes \texttt{CLAUDE.md} at the repository
root, providing system-level instructions for AI coding
assistants---architectural principles, testing requirements, and naming
conventions that apply globally. This creates a three-tier documentation
architecture: per-directory \texttt{AGENTS.md} for local context, root
\texttt{README.md} and \texttt{CLAUDE.md} for global constraints, and
\texttt{README.md} for human navigation.
\end{frame}

\begin{frame}[fragile,allowframebreaks]{Agentic Skill Architecture}
\protect\phantomsection\label{agentic-skill-architecture}
The \hyperlink{documentation-duality-and-ai-collaboration}{Documentation
Duality} standard addresses human and AI navigation at the directory
level. A complementary layer operates at the \emph{module} level: every
infrastructure subpackage carries two additional machine-readable files
that transform it from a passive code library into an active,
protocol-aligned skill endpoint.

\begin{block}{The Three-Tier Skill Protocol}
\protect\phantomsection\label{the-three-tier-skill-protocol}
\texttt{template/} implements a progression of agent-facing
documentation, escalating in specificity from global constraints to
module-level API contracts:

{\def\LTcaptype{none} % do not increment counter
\begin{longtable}[]{@{}
  >{\raggedright\arraybackslash}p{(\linewidth - 6\tabcolsep) * \real{0.2143}}
  >{\raggedright\arraybackslash}p{(\linewidth - 6\tabcolsep) * \real{0.2143}}
  >{\raggedright\arraybackslash}p{(\linewidth - 6\tabcolsep) * \real{0.2500}}
  >{\raggedright\arraybackslash}p{(\linewidth - 6\tabcolsep) * \real{0.3214}}@{}}
\toprule\noalign{}
\begin{minipage}[b]{\linewidth}\raggedright
Tier
\end{minipage} & \begin{minipage}[b]{\linewidth}\raggedright
File
\end{minipage} & \begin{minipage}[b]{\linewidth}\raggedright
Scope
\end{minipage} & \begin{minipage}[b]{\linewidth}\raggedright
Purpose
\end{minipage} \\
\midrule\noalign{}
\endhead
1 --- System & \texttt{README.md} & Repository root & Global
architectural principles, Zero-Mock policy, naming conventions \\
2 --- Structure & \texttt{AGENTS.md} & Every directory & Local file
inventories, API surfaces, integration patterns, architectural
constraints \\
3 --- Skill & \texttt{SKILL.md} & Every infrastructure module &
Machine-parseable skill descriptor: module name, description, key
imports, usage pattern \\
\bottomrule\noalign{}
\end{longtable}
}

Tier 1 and Tier 2 have direct analogues in the prompt-engineering
literature: system prompts and retrieval-augmented context
\breakseq{[@lau2025aicoding]}. Tier 3 is novel. The \texttt{SKILL.md} files
follow a YAML frontmatter + Markdown instruction format precisely
aligned with the tool-descriptor schemas of the Model Context Protocol
\breakseq{[@anthropic2024mcp]}. The following is the exact frontmatter from
\breaktt{infrastructure/rendering/SKILL.md}:

\begin{Shaded}
\begin{Highlighting}[]
\PreprocessorTok{{-}{-}{-}}
\FunctionTok{name}\KeywordTok{:}\AttributeTok{ rendering}
\FunctionTok{description}\KeywordTok{: }\CharTok{\textgreater{}}
\NormalTok{  Multi{-}format output generation (PDF, HTML, slides).}
\NormalTok{  Use for: Pandoc/XeLaTeX rendering, RenderManager, slide deck generation.}
\NormalTok{  Key imports: RenderManager, RenderingConfig from infrastructure.rendering}
\PreprocessorTok{{-}{-}{-}}
\end{Highlighting}
\end{Shaded}

An MCP client reading this block immediately knows the module name, its
natural-language activation condition (``use for''), and which Python
symbols to import. No source-code inspection is required. This is the
practical implementation of Toolformer-style self-documented tools
\breakseq{[@schick2023toolformer]}---rather than a language model learning tool
APIs from training data, the APIs are encoded directly in
version-controlled, co-located skill files that evolve with the
codebase.
\end{block}

\begin{block}{Module Skill Coverage}
\protect\phantomsection\label{module-skill-coverage}
Each infrastructure subdirectory surfaced by
\breaktt{discover\_infrastructure\_modules()} carries paired
\texttt{README.md} + \texttt{AGENTS.md}; agent-facing \texttt{SKILL.md}
manifests exist wherever teams enable Cursor / PAI manifests
(regenerated via \breaktt{python\ -m\ infrastructure.skills}). Root-level
\texttt{PAI.md} summarizes cross-package obligations.

Promotion policy: new Layer‑1 directories must ship human +
machine-readable docs (\texttt{README.md}, \texttt{AGENTS.md})
immediately; Tier‑3 SKILL assets follow once APIs stabilize.
\end{block}

\begin{block}{MCP Server Mapping}
\protect\phantomsection\label{mcp-server-mapping}
The mapping from \texttt{SKILL.md} descriptors to MCP server endpoints
is intentional but not yet automated; it represents the principal next
integration step. In the MCP architecture \breakseq{[@anthropic2024mcp]}, a
server exposes three primitive types: \textbf{Tools} (executable
functions), \textbf{Resources} (data the model can read), and
\textbf{Prompts} (structured query templates). Each
\texttt{infrastructure} module maps naturally onto this taxonomy:

\begin{itemize}
\tightlist
\item
  \breaktt{infrastructure.llm} → MCP \textbf{Tool} (\texttt{query},
  \texttt{apply\_template}) + MCP \textbf{Prompt} (research prompt
  templates)
\item
  \breaktt{infrastructure.rendering} → MCP \textbf{Tool}
  (\texttt{render\_pdf}, \texttt{render\_html}) + MCP \textbf{Resource}
  (rendered PDFs as URI-addressable resources)
\item
  \breaktt{infrastructure.validation} → MCP \textbf{Tool}
  (\breaktt{validate\_pdf\_rendering}, \breaktt{validate\_markdown})
\item
  \breaktt{infrastructure.publishing} → MCP \textbf{Tool}
  (\breaktt{publish\_to\_zenodo}, \breaktt{generate\_citation\_bibtex}) +
  MCP \textbf{Resource} (DOI registry)
\item
  \breaktt{infrastructure.steganography} → MCP \textbf{Tool}
  (\breaktt{SteganographyProcessor.process}) + MCP \textbf{Resource}
  (hash manifests)
\item
  \breaktt{infrastructure.search} · \breaktt{infrastructure.reference} →
  MCP \textbf{Tool} wrappers over literature retrieval + BibTeX handling
  + MCP \textbf{Resource} exports for corpus JSON / \texttt{.bib}
\end{itemize}

An agent orchestrating a full research pipeline could, in principle,
compose these MCP tools to reproduce the declarative DAG
programmatically---discovering capabilities via \texttt{SKILL.md}
frontmatter, executing them via MCP protocol calls, and consuming their
outputs as Resources. The \texttt{SKILL.md} files parallel Voyager's
skill library \breakseq{[@wang2023voyager]}---Voyager's agent accumulates a
growing library of executable Minecraft skills represented as JavaScript
functions; \texttt{template/}'s agent accumulates a curated library of
research pipeline skills represented as YAML-frontmattered
\texttt{SKILL.md} files. In both cases, the skill representation is
machine-readable, version-controlled, and self-describing. Wang et al.'s
LLM agent survey \breakseq{[@wang2024llmagents]} identifies tool use, planning,
and memory as the three fundamental capabilities of autonomous agents;
Yao et al.'s ReAct framework \breakseq{[@yao2023react]} demonstrates that
interleaving reasoning traces with tool actions dramatically improves
agent reliability in interactive settings. The \texttt{template/} skill
architecture maps cleanly onto these three capabilities: the
\texttt{SKILL.md} descriptors supply the tool-use layer, the declarative
DAG of \texttt{12} \texttt{pipeline.yaml} stages (a default full run
executes \texttt{10}) supplies the planning scaffold, and the
per-criterion checkpoint system supplies the memory layer.
\end{block}
\end{frame}

\end{document}
